\subsection{Synthesis}

The two halves bound the operating envelope of marker terminal docking from opposite sides. Part~A bounds the dynamics. With the full perception pipeline in the loop, the vehicle captures a swaying dock cleanly down to 8~s periods at 0.1~m amplitude, limited by its own control bandwidth (the 240~ms plant lag), and delivers insertions that a capture funnel of 0.15~m radius and 0.1~m/s tolerance can accept. Part~B bounds the sensing, where a marker is reliably detectable only above $\sim$27~px of apparent size and below $\tau \approx 2$ of optical depth. In the measured water, this sets the 5~m acquisition marker at 182~mm (154 to 207~mm), and it rules out reliable acquisition at that range in substantially murkier water regardless of the marker size.

The intersection of the two envelopes is the design statement for the dock. The 200~mm-class markers assumed throughout Part~A are validated by Part~B as sufficient for 5~m acquisition in pool-clarity water. The per-size fusion pipeline and the size-graded layout of Part~A (Section~\ref{sec:method-layout}) consume the same size-dependent envelope that Part~B measured, since the small markers take over as the large ones overflow the FoV at close range, and the 5~cm blackout sets the seat criterion. Beyond the measured envelope, at longer standoff or higher turbidity, the marker channel provides no signal. Terminal docking must then be handed a vehicle already inside the envelope by a far-field cue such as an active light beacon, which is the acquisition/search phase this thesis scoped out.
